% Lösungen Einheit 1 von 4 – Wurzelziehen und Lösbarkeit (zusammenbau v0.1)
\einheitenkopf{Einheit 1 von 4 · Wurzelziehen und Lösbarkeit}

\erg{11}{a) quadratisch – x kommt im Quadrat vor}
%% TODO Lösung der Erklärzeile 12g) fehlt in der Bank
\erg{12}{a) zwei Lösungen – rechts positiv \quad b) $x_1 = 9$, $x_2 = -9$ \quad c) $x_1 = 5$, $x_2 = -5$ \quad d) $x_1 = 10$, $x_2 = -10$ \quad e) $x_1 = 13$, $x_2 = -13$ \quad f) $x_1 = 14$, $x_2 = -14$ \quad g) \ldots \quad h) Seitenlänge $\sqrt{64} = 8$ m \quad i) $x_1 = \sqrt{13} \approx 3{,}61$, $x_2 = -\sqrt{13} \approx -3{,}61$ \quad j) keine Lösung: $x^2 = -100$ ist unmöglich, weil ein Quadrat nie negativ ist \quad k) $x^2 = 25$; $x_1 = 5$, $x_2 = -5$ \quad l) $r^2 = 200 : (\pi \cdot 8) \approx 7{,}96$; $r \approx 2{,}82$ dm \quad m) $x^2 = 16$; $x_1 = 4$, $x_2 = -4$ \quad n) $x - 6 = 4$ oder $x - 6 = -4$; $x_1 = 10$, $x_2 = 2$ \quad o) $x + 10 = 5$ oder $x + 10 = -5$; $x_1 = -5$, $x_2 = -15$ \quad p) genau eine Lösung: $x = 9$ \quad q) $(-7 + 2)^2 = (-5)^2 = 25$ (wA): ja \quad r) $2$ – zwei Lösungen, der Scheitel liegt unter der Geraden \quad s) z. B. $x^2 = -3$ – jede negative Zahl rechts \quad t) A wahr: nur $x = 6$; B falsch: $x + 9 = 2$ oder $x + 9 = -2$, also $x_1 = -7$, $x_2 = -11$; ohne Lösung z. B. $x^2 = -7$}
\erg{13}{a) Die negative Lösung fehlt. Richtig: $x_1 = 15$, $x_2 = -15$ \quad b) $\sqrt{10}$ und $-\sqrt{10}$ ergeben quadriert beide $10$, denn minus mal minus ist plus \quad c) $x_1 = 2$, $x_2 = -2$ \quad d) $t^2 = 80 : 5 = 16$; $t = 4$ s (die negative Lösung ist keine Zeit)}
